\chapter{Conclusioni}
\label{chap:conclusioni}

Questo capitolo conclude il lavoro con un consuntivo finale rispetto agli obiettivi del Piano di lavoro, una valutazione retrospettiva dell'esperienza di stage -- sul piano del raggiungimento degli obiettivi, della crescita professionale personale e del confronto fra le competenze richieste e quelle fornite dal corso di studi -- e una discussione dei possibili sviluppi futuri del lavoro.

\section{Consuntivo finale}
\label{sec:consuntivo-finale}

\todoplaceholder{consuntivo finale rispetto agli obiettivi obbligatori, desiderabili e opzionali del Piano di lavoro, e ai deliverable prodotti.}

% TODO: consuntivo finale del lavoro svolto, distinguendo esplicitamente il
% grado di raggiungimento degli obiettivi obbligatori (analisi e
% classificazione dei criteri WCAG, prototipo di Accessibility Checker,
% metodologia di validazione, validazione sperimentale, documentazione),
% desiderabili (miglioramento del prompt engineering, confronto tra LLM,
% studio di Pi Agent e OpenCode, implementazione di almeno una skill
% agentica, confronto fra approccio a prompting e approccio agentico) e
% opzionali (skill agentiche aggiuntive, confronto tra modelli più ampio,
% benchmark più esteso); elenco dei deliverable effettivamente prodotti,
% rispetto a quelli previsti nel Piano di lavoro.


\section{Raggiungimento degli obiettivi}
\label{sec:raggiungimento-obiettivi}

\todoplaceholder{grado di soddisfacimento degli obiettivi, su base fattuale.}

% TODO: valutazione, basata su dati di fatto e non su opinioni personali, del
% grado di soddisfacimento degli obiettivi indicati nel capitolo~\ref{chap:motivazioni-obiettivi},
% sia dal punto di vista dello studente sia da quello dell'organizzazione ospitante.


\section{Crescita professionale}
\label{sec:crescita-professionale}

\todoplaceholder{conoscenze e abilità acquisite durante lo stage.}

% TODO: riflessione su cosa lo stage ha effettivamente portato in termini di
% conoscenze e abilità acquisite, con riferimento concreto alle attività
% descritte nei capitoli~\ref{chap:metodologia-selezione}--\ref{chap:estensione-agentica}.


\section{Confronto tra competenze richieste e percorso di studi}
\label{sec:confronto-competenze}

\todoplaceholder{distanza tra competenze richieste e percorso di studi.}

% TODO: distanza tra le competenze richieste a inizio stage per svolgere il
% lavoro e quelle fornite dal corso di studi, con eventuale segnalazione di
% lacune riscontrate.


\section{Sviluppi futuri}
\label{sec:sviluppi-futuri}

\todoplaceholder{possibili sviluppi futuri del lavoro.}

% TODO: possibili sviluppi futuri del sistema, distinguendo quelli relativi
% al blocco principale (ad esempio ampliamento dei criteri core, benchmark
% più esteso) da quelli relativi all'estensione agentica (ad esempio un
% workflow agentico più esteso, ulteriore validazione), coerentemente con gli
% obiettivi opzionali indicati nel Piano di lavoro.
